\documentclass[12pt,a4paper]{article}
\usepackage[utf8]{inputenc}
\usepackage[T1]{fontenc}
\usepackage{amsmath, amsfonts, amssymb}
\usepackage{geometry}
\usepackage{lmodern}
\usepackage[french]{babel}
\usepackage[colorlinks=true, linkcolor=blue]{hyperref}

\geometry{hmargin=2cm,vmargin=1.5cm}
\setlength{\parindent}{0pt}
\setlength{\parskip}{5pt}

\begin{document}

%page 1 
\begin{center}
    \rule{\linewidth}{0.5pt} \\
    \Large \textbf{Equation de Korteweg-de Vries} \\
    \rule{\linewidth}{0.5pt} \\
    \vspace{0.5cm}
    \[ \frac{\partial u(x,t)}{\partial t}+\epsilon u(x,t)\frac{\partial u(x,t)}{\partial x}+\mu\frac{\partial^{3}u(x,t)}{\partial x^{3}}=0 \]
\end{center}


\section{Construction du Schéma par différence finie}


On discrétise à la fois la variable x et la variable t. On définit pour ça un maillage ;  $u_i^n$ est la solution approchée au point $x_i = i\Delta x$ et à l'instant $t_n = n\Delta t$. 

\subsection{En espace}
Pour la dérivée première en espace on utilise la formule de différence finie centrée, qui garantit une précision d'ordre 2 : $\frac{\partial u}{\partial x}(x_i) \approx \frac{u_{i+1} - u_{i-1}}{2\Delta x}$. 

Pour la dérivée troisième, on utilise les formules de Taylor Young pour construire une approximation à l'ordre 3 (\hyperref[annexe:demo1]{(1)} démonstration). 
$\frac{\partial^3 u}{\partial x^3}(x_i) \approx \frac{u_{i+2} - 2u_{i+1} + 2u_{i-1} - u_{i-2}}{2\Delta x^3}$

On a donc $\frac{d u_i}{dt} = - \epsilon u_i \frac{u_{i+1} - u_{i-1}}{2\Delta x} - \mu \frac{u_{i+2} - 2u_{i+1} + 2u_{i-1} - u_{i-2}}{2\Delta x^3}$

\subsection{En temps}

Pour le temps : 

-est ce qu'on utilise une formule de DF centrée classique, mais du coup on a besoin du temps futur donc ça complique l'histoire (je sais même pas comment résoudre ça (systèmes de matrices probablement + conditions aux bords ?)

-On peut utiliser un df progressive $\frac{\partial u}{\partial t}(x_i, t_n) \approx \frac{u_i^{n+1} - u_i^n}{\Delta t}$, ensuite on isole $u_i^{n+1}$ d'un coté de l'eq car c'est ce qu'on cherche (le point actuel au temps $t+\Delta t$)

-On utilise runge kutta, il faut s'y replonger

-autre ? 


\subsection{Conditions limites}

Dans tout les problèmes qu'on a connu on avait des CL et un domaine fermé mais est ce qu'on veut un dmaine fermé ici ? En soit on est pas obliger. On cherche à étudier le mouvement de la vague sans limites. 
Et du coup si on en a pas on pourrai aller piocher les points hors domaines (ex $u_{i-2}$ en $i=1$) à la fin de la matrice, ce qui ferait donc une "boucle". 

\subsection{Étude du schéma}

Consistance, stabilité et convergence à vérifier ou bien ? 

\newpage

\phantomsection
\section*{Annexe}
\phantomsection\label{annexe:demo1}
\subsubsection*{\hyperref[main:demo1]{(1)} Démonstration du schéma aux différences finies centrées pour la dérivée d'ordre 3}

L'objectif est d'établir l'approximation par différences finies de la dérivée troisième $u^{(3)}(x)$.

Supposons la solution $u$ de classe $\mathcal{C}^5$.

Considérons la formule de Taylor-Lagrange à l'ordre 5 appliquée à la solution $u$ entre $x$ et $x+h$. Il existe $\theta_1 \in ]0,1[$ tel que:
\begin{equation}
u(x+h) = u(x) + h u'(x) + \frac{1}{2}h^2 u''(x) + \frac{1}{6}h^3 u^{(3)}(x) + \frac{1}{24}h^4 u^{(4)}(x) + \frac{1}{120}h^5 u^{(5)}(x+\theta_1 h). \label{eq:t+}
\end{equation}

Par ailleurs, d'après la formule de Taylor-Lagrange à l'ordre 5 appliquée à la solution $u$ entre $x$ et $x-h$, il existe $\theta_2 \in ]-1,0[$ tel que:
\begin{equation}
u(x-h) = u(x) - h u'(x) + \frac{1}{2}h^2 u''(x) - \frac{1}{6}h^3 u^{(3)}(x) + \frac{1}{24}h^4 u^{(4)}(x) - \frac{1}{120}h^5 u^{(5)}(x+\theta_2 h). \label{eq:t-}
\end{equation}

On soustrait $(\ref{eq:t+})  et  (\ref{eq:t-})$ :
\begin{equation}
u(x+h) - u(x-h) = 2h u'(x) + \frac{1}{3}h^3 u^{(3)}(x) + \frac{h^5}{120}\left( u^{(5)}(x+\theta_1 h) + u^{(5)}(x+\theta_2 h) \right). \label{eq:diff_h}
\end{equation}


On cherche à isoler $u^{(3)}(x)$, l'équation contient encore $u'(x)$, on élargit aux points $x+2h$ et $x-2h$. 

Il existe $\theta_3 \in ]0,1[$ et $\theta_4 \in ]-1,0[$ tels que :
\begin{equation}
u(x+2h) = u(x) + 2h u'(x) + \frac{4h^2}{2} u''(x) + \frac{8h^3}{6} u^{(3)}(x) + \frac{16h^4}{24} u^{(4)}(x) + \frac{32h^5}{120} u^{(5)}(x+2\theta_3 h),
\end{equation}
\begin{equation}
u(x-2h) = u(x) - 2h u'(x) + \frac{4h^2}{2} u''(x) - \frac{8h^3}{6} u^{(3)}(x) + \frac{16h^4}{24} u^{(4)}(x) - \frac{32h^5}{120} u^{(5)}(x+2\theta_4 h).
\end{equation}

On soustrait et on obtient :
\begin{equation}
u(x+2h) - u(x-2h) = 4h u'(x) + \frac{8}{3}h^3 u^{(3)}(x) + \frac{32h^5}{120}\left( u^{(5)}(x+2\theta_3 h) + u^{(5)}(x+2\theta_4 h) \right). \label{eq:diff_2h}
\end{equation}

Nous avons un système de deux équations, (\ref{eq:diff_h}) et (\ref{eq:diff_2h}). On cherche à éliminer le terme en $u'(x)$ : 
$$(\ref{eq:diff_2h}) - 2 \times (\ref{eq:diff_h})$$

Ce qui donne  : 

\begin{align*}
u(x+2h) - 2u(x+h) + 2u(x-h) - u(x-2h) &= 2h^3 u^{(3)}(x) \\
& + \frac{4h^5}{15}\left( u^{(5)}(x+2\theta_3 h) + u^{(5)}(x+2\theta_4 h) \right) \\
&\quad - \frac{h^5}{60}\left( u^{(5)}(x+\theta_1 h) + u^{(5)}(x+\theta_2 h) \right).
\end{align*}

En isolant $u^{(3)}(x)$ on a : 

\begin{equation}
u^{(3)}(x) = \frac{u(x+2h) - 2u(x+h) + 2u(x-h) - u(x-2h)}{2h^3} + \epsilon,
\end{equation}
avec $\epsilon$ l'erreur de consistance 
\begin{equation}
\epsilon = - \frac{h^2}{2} \left[ \frac{4}{15}\left( u^{(5)}(x+2\theta_3 h) + u^{(5)}(x+2\theta_4 h) \right) - \frac{1}{60}\left( u^{(5)}(x+\theta_1 h) + u^{(5)}(x+\theta_2 h) \right) \right]. \label{eq:erreur_exacte}
\end{equation}

\vspace{5cm}
\textbf{(à faire potentiellement dans une autre section)}

Majoration de l'erreur de consistance 

Notons $M_5 = \sup_{t \in \mathcal{D}} |u^{(5)}(t)|$.

\begin{align*}
|\epsilon| &\leq \frac{h^2}{2} \left[ \frac{4}{15} \left( |u^{(5)}(x+2\theta_3 h)| + |u^{(5)}(x+2\theta_4 h)| \right) + \frac{1}{60} \left( |u^{(5)}(x+\theta_1 h)| + |u^{(5)}(x+\theta_2 h)| \right) \right] \\
&\leq \frac{h^2}{2} \left[ \frac{4}{15} (2 M_5) + \frac{1}{60} (2 M_5) \right] \\
&\leq \frac{h^2}{2} \left[ \frac{8}{15} + \frac{2}{60} \right] M_5.
\end{align*}


Ainsi, 
\begin{equation}
|\epsilon| \leq \frac{h^2}{2} \times \frac{17}{30} M_5 = \frac{17}{60} h^2 M_5.
\end{equation}

Le schéma aux différences finies centrées pour la dérivée troisième est consistant d'ordre 2. 

\end{document}
